\section{综合案例：课程知识问答助手的系统设计}
\label{sec:case}

\begin{guide}
\noindent 假设要为一门课程搭建知识问答助手，它既要能解释课程概念，又要给出讲义出处。该怎么搭、需要什么硬件、速度够不够？本节按自上而下的顺序回答：业务需求决定服务指标，服务指标决定模型与推理引擎，模型与负载决定显存和带宽需求，最后才确定硬件。每一步都先用本章的估算公式判断方案是否可行，再用实测验证并迭代。
\end{guide}

本节把前面各节的知识串联起来，完整演示一次系统设计。设计的顺序与本章的组织思路一致：\emph{先明确需求，再选择方案，最后评估效果}（图\,\ref{fig:case-flow}）。

\begin{figure}[htbp]
\centering
\begin{tikzpicture}[fig, x=1mm, y=1mm]
  % 两栏对齐：左栏为设计环节，右栏为本案例取值；整体以 x=0 为中心
  \node[note, font=\figfont\footnotesize\bfseries] at (-16,9) {设计环节};
  \node[note, font=\figfont\footnotesize\bfseries] at (32,9) {本案例中的取值};
  \foreach \i/\t/\c in {
      0/业务需求/{200 名学生，高峰 30 路并发\\回答注明出处，数据不外传},
      1/服务指标/{P95 TTFT ≤ 2 s，输出 ≥ 15 token/s\\引用可核对},
      2/模型与推理引擎/{RAG 工作流 + 8B 开源模型\\vLLM 类推理引擎},
      3/显存与带宽估算/{4 位权重约 5 GB\\KV 缓存每请求约 0.5 GiB},
      4/硬件与网络/{单张 24 GB 显卡\\显存带宽约 1 TB/s}} {
    \ifnum\i=3
      \node[blka, minimum width=32mm, minimum height=10mm] (s\i) at (-16,-\i*14) {\t};
    \else
      \node[blk, minimum width=32mm, minimum height=10mm] (s\i) at (-16,-\i*14) {\t};
    \fi
    \node[blkw, draw=rulegray, minimum width=50mm, minimum height=10mm, font=\figfont\footnotesize, align=left, text width=46mm] (v\i) at (32,-\i*14) {\c};
    \draw[densely dotted, rulegray, line width=0.6pt] (s\i.east) -- (v\i.west);
  }
  \foreach \i [evaluate=\i as \j using int(\i+1)] in {0,...,3}
    \draw[flow] (s\i.south) -- (s\j.north);
  % 实测反馈回路（左侧）
  \node[blkw, minimum width=15mm, minimum height=10mm, font=\figfont\footnotesize, align=center] (m) at (-50,-35) {实测\\验证};
  \draw[->, dashed, line width=0.6pt] (s4.west) -| (m.south);
  \draw[->, dashed, line width=0.6pt] (m.north) |- (s1.west);
  % 居中图例
  \begin{scope}[shift={(0,-67)}]
    \draw[flow] (-40,0) -- (-34,0); \node[note, anchor=west, font=\figfont\scriptsize] at (-33,0) {设计方向：自上而下};
    \draw[->, dashed, line width=0.6pt] (2,0) -- (8,0); \node[note, anchor=west, font=\figfont\scriptsize] at (9,0) {用测量结果修正指标与方案};
  \end{scope}
\end{tikzpicture}
\caption{自上而下的系统设计流程及其在课程知识问答助手中的取值。高亮的估算环节决定了方案能否成立}
\label{fig:case-flow}
\end{figure}


\subsection{第一步：明确需求与服务指标}

某高校希望为一门课程开发问答助手，基本情况如下：
\begin{itemize}
  \item \emph{用户与负载}：选课学生约 200 人，考试前的高峰时段约有 30 人同时提问；
  \item \emph{数据}：课件和讲义共约 2000 页，每学期更新；
  \item \emph{功能}：回答需注明出自哪份讲义的哪一页，便于学生核对；
  \item \emph{约束}：课程资料不得上传到外部服务，预算为一台配备单张 \SI{24}{GB} 消费级显卡的服务器。
\end{itemize}

将需求转化为可度量的服务指标：P95 的 TTFT 不超过 \SI{2}{s}；输出速度不低于每秒 15 个 token（明显快于阅读速度）；在 30 路并发下服务保持稳定；引用页码的正确率需通过人工抽查评估。

\subsection{第二步：自上而下选择方案}

\paragraph{应用层：先 RAG，不微调} 需求的核心是“知识准确且可溯源”，而且资料每学期更新，这正是 RAG 擅长的任务（见第\,\ref{sec:tools}\,节的辨析）。因此不需要训练或微调模型。处理流程固定为：解析讲义并保存课程、讲次、页码等元数据；对问题同时做向量检索和关键词检索；把检索到的片段与问题一起交给模型；生成带引用的回答。由于步骤固定，采用工作流而非智能体，便于测试和排错。2000 页资料切分后约为数万个片段，规模很小，单机向量库即可满足需求。

\paragraph{模型层：选择规模与精度} 数据不能外传，因此在本地部署开源权重模型。考虑到显存只有 \SI{24}{GB}，选择约 8B 参数的指令模型。由式\,\eqref{eq:weight-mem} 估算：BF16 权重约 \SI{16}{GB}，4 位量化后约 \SI{5}{GB}（含元数据）。

\paragraph{服务层：选择推理引擎} 高峰期有 30 路并发，需要连续批处理和高效的 KV 缓存管理，因此选用 vLLM 一类的推理引擎，而不是以单用户体验为主的本地运行工具。

\subsection{第三步：估算效果，检查可行性}

\paragraph{显存是否够用} 每个请求的上下文按 4096 个 token 计（问题、检索片段与回答之和），由例\,\ref{ex:kv} 估算，每个请求的 KV 缓存约 \SI{0.5}{GiB}。扣除运行时开销（按约 \SI{2}{GB} 计）后：
\begin{itemize}
  \item BF16 权重：剩余约 \SI{6}{GB}，仅能容纳约 11 个请求，\emph{无法满足} 30 路并发；
  \item 4 位量化权重：剩余约 \SI{17}{GB}，可容纳约 31 个请求，\emph{基本满足}需求，但余量很小。
\end{itemize}
因此选择 4 位量化。若实测发现回答质量明显下降，可以考虑缩短检索上下文、对 KV 缓存也进行量化，或者增加一张显卡。

\paragraph{速度是否够快} 该显卡的显存带宽约 \SI{1}{TB/s}。按第\,\ref{subsec:three-quantities}\,小节的方法，4 位权重下单请求解码速度的理论上限约为 $1000/5=200$ token/s。即便实际只能达到上限的几分之一，且多个请求共享带宽，也有较大把握满足每秒 15 个 token 的目标。预填充阶段，3000 个输入 token 约需 $2\times8\times10^{9}\times3000\approx5\times10^{13}$ 次浮点运算，按有效算力约 $10^{14}$\,FLOP/s 计，耗时约 \SI{0.5}{s}，满足 TTFT 目标；高峰期的排队时间则需要实测确认。

\paragraph{瓶颈可能在哪里} 以上估算表明，在这个场景中，主要约束是\emph{显存容量}（决定并发上限），而不是峰值算力。如果实测发现响应慢的主要原因在检索环节，就应优化索引和重排序，而不是更换显卡。

\subsection{第四步：实测验证}

估算只能判断方案大致是否可行，最终必须以实测为准。建议开展以下三项实验：
\begin{enumerate}
  \item \emph{单请求基线}：固定输出长度，分别使用 512、2048、4096 个 token 的输入，记录 TTFT 与输出速度，检验预填充时间随输入长度的变化；
  \item \emph{并发压力测试}：保持模型、硬件和提示词不变，把并发数从 1 逐步增加到 40，绘制总吞吐与 P95 延迟随并发的变化曲线，找出 KV 缓存耗尽、开始排队的拐点；
  \item \emph{检索效果检验}：向索引中加入或移除一份带页码的讲义，观察回答内容与引用是否随之变化，并人工抽查引用的正确率。
\end{enumerate}
实验报告应记录硬件型号、驱动与引擎版本、模型与量化方式、输入输出长度分布和并发设置，确保结果可以复现。

\begin{keypoints}
这个案例体现了本章的核心思路：\emph{业务需求决定服务指标，服务指标决定模型与推理引擎，模型与负载决定存储与带宽需求，存储、带宽与并发需求最终决定硬件和网络}。只记住工具的名称，无法完成这样的设计；掌握各层的作用与估算方法，才能在动手之前判断方案是否可行。
\end{keypoints}
